\documentclass[10pt,a4paper]{amsart}
\usepackage[margin=19mm]{geometry}
\usepackage{amsmath,amssymb,mathtools}
\usepackage[scr=rsfs,cal=euler]{mathalfa}
\usepackage{xfrac}
\usepackage[unicode,hidelinks,bookmarks=false]{hyperref}
\newcommand{\ie}{i.e.\ }
\newcommand{\cf}{cf.\ }
\newcommand{\resp}{respectively\ }
\newcommand{\Subsection}{\subsection}
\providecommand{\nobreakdash}{\nobreak\hbox{-}}
\setlength{\emergencystretch}{2em}
\multlinegap0pt
\usepackage{bm}
\usepackage{bbm}
\usepackage{dsfont}
\usepackage{xspace}
\usepackage{cancel}
\usepackage{mathtools}
\usepackage{amsthm}
\makeatletter
\@ifundefined{theorem}{\newtheorem{theorem}{Theorem}}{}
\makeatother
\newcommand{\R}{\mathbb{R}}
\newcommand{\sym}{\mathrm{sym}}
\title[The Monge-Ampere system: convex integration in arbitrary dim]{The Monge-Ampere system: convex integration in arbitrary dimension and codimension}
\author{Marta Lewicka}
\date{Source: arXiv:2210.04363v4 (CC BY 4.0)}
\begin{document}
\maketitle

\noindent\emph{Source and licence.} The Monge-Ampere system: convex integration in arbitrary dimension and codimension. Version arXiv:2210.04363v4 (CC BY 4.0), distributed under the Creative Commons Attribution 4.0 International licence, as recorded in the arXiv licence metadata for this exact version and shown on the abstract page https://arxiv.org/abs/2210.04363v4. Full rights notice: licenses/arxiv-2210.04363-CC-BY-4.0.txt.\par\smallskip

\noindent\textit{Editorial scope.} Counted unit: the Theorem whose source label is \texttt{th\_NashKuiHol} in the article (source0.tex lines 1800--1825, source label \texttt{th\_NashKuiHol}), delivered together with the article's own complete original proof of it (source0.tex lines 1826--1970, one \texttt{proof} environment of 145 lines). No mathematics is written, repaired or re-derived: statement and proof are the article's own text line for line. Required context retained uncounted: thm\_stage (lines 305--329); Abound2 (lines 318--326). Every definition, display and label of those lines is retained unchanged. Omitted from this package and not redistributed: 1--304, 327--1799, 1971--2677. Nothing inside a retained excerpt is omitted or altered: every retained body is delivered exactly as the article's lines are written, so no omission receipt of any kind accompanies this package. Citations: the retained excerpts carry no citation command at all, so no bibliography entry of the article is transcribed and no reference is redistributed. Reference adaptation: none was needed and none was made; every reference inside the delivered unit resolves inside it, so no unresolved reference or citation is produced. Printed slips kept literal: no printed word, symbol, hypothesis or number of the counted statement or of its proof was altered. Layout and class: the article's own class is \texttt{amsart} with packages including amsmath, amsfonts, amssymb, esint, amsthm, nicefrac, bbm, dsfont, epsfig, graphicx; the wrapper is the lane's pinned \texttt{amsart} editorial wrapper at 10pt on A4, which restates the theorem-like environments the retained text needs and adds the article's own macro definitions that the retained text needs (\texttt{\textbackslash R}, \texttt{\textbackslash sym}), with extra wrapper packages (bm, bbm, dsfont, xspace, cancel, mathtools). The delivered numbering is the unit's own. Assets: no retained span contains a figure environment, a graphics-inclusion command, a PDF-page inclusion, a table or a TikZ picture; no image file is copied into this package, the package declares no asset dependency and no image file is redistributed with it. Licence: version arXiv:2210.04363v4 of this article is distributed under the Creative Commons Attribution 4.0 International licence, as recorded in the arXiv licence metadata for this exact version and shown on the abstract page https://arxiv.org/abs/2210.04363v4. A full rights notice is delivered at licenses/arxiv-2210.04363-CC-BY-4.0.txt. Characters: every retained excerpt is pure ASCII, so the delivered engine, XeLaTeX with its default Latin Modern text font, reports no missing character.
\begin{theorem}\label{thm_stage}
Let the vector fields $v\in\mathcal{C}^2(\bar \omega,\R^k)$, $w\in\mathcal{C}^2(\bar\omega,\R^d)$ and
the matrix field $A\in\mathcal{C}^{0,\beta}(\bar\omega,\R^{d\times d}_\sym)$ be given on an open,
bounded domain $\omega\subset\R^d$. Assume that:
$$\mathcal{D}=A -\big(\frac{1}{2}(\nabla v)^T\nabla v + \sym\nabla
w\big) \quad \mbox{ satisfies } \; 0<\|\mathcal{D}\|_0\leq 1.$$
Fix two constants $M, \sigma$ such that:
$$M\geq\max\{\|v\|_2, \|w\|_2, 1\} \quad\mbox{ and }\quad \sigma\geq 1.$$
Then, there exist $\tilde v\in\mathcal{C}^2(\bar \omega,\R^k)$ and
$\tilde w\in\mathcal{C}^2(\bar\omega,\R^d)$ such that, denoting:
$$\tilde{\mathcal{D}}=A -\big(\frac{1}{2}(\nabla \tilde v)^T\nabla \tilde v + \sym\nabla \tilde w\big), $$
the following holds: 
%\begin{equation}\label{stage_est}
\begin{align*}
& \|\tilde v - v\|_1\leq C\|\mathcal{D}\|_0^{1/2}, \quad \|\tilde w -
w\|_1\leq C\|\mathcal{D}\|_0^{1/2}(1+\|\nabla v\|_0),
\tag*{(\theequation)$_1$}\refstepcounter{equation} \label{Abound1}\vspace{1mm}\\
& \|\nabla^2\tilde v\|_0\leq CM\sigma^{d_*/k},\quad \|\nabla^2\tilde w\|_0\leq CM\sigma^{d_*/k}(1+\|\nabla v\|_0),
\tag*{(\theequation)$_2$}\label{Abound2} \\ 
& \|\tilde{\mathcal{D}}\|_0\leq C\Big(\frac{\|A\|_{0,\beta}}{M^\beta} \|\mathcal{D}\|_0^{\beta/2} +
\frac{\|\mathcal{D}\|_0}{\sigma}\Big), \tag*{(\theequation)$_3$} \label{Abound3}
\end{align*}
%\end{equation}
where $d_*=d(d+1)/2$ and where the constants $C$ depend only on $d, k$ and $\omega$.
\end{theorem}

\begin{theorem}\label{th_NashKuiHol}
Let $\omega\subset\R^d$ be an open, bounded domain 
and let $k\geq 1$ and $\gamma>0$ be such that the statement of Theorem \ref{thm_stage}
holds true with $\gamma$ replacing the exponent $d_*/k$ in
\ref{Abound2}, provided that $\sigma>\sigma_0$ where $\sigma_0>1$
depends only on $\omega$ and $\gamma$. Then we have the following.  
For every $v\in\mathcal{C}^2(\bar\omega,\R^k)$, $w\in\mathcal{C}^2(\bar\omega,\R^d)$,
$A\in\mathcal{C}^{0,\beta}(\bar\omega, \R^{d\times d}_\sym)$, such that:
$$\mathcal{D}=A-\big(\frac{1}{2}(\nabla v)^T\nabla v + \sym\nabla
w\big) \quad\mbox{ satisfies } \quad 0<\|\mathcal{D}\|_0\leq 1,$$
and for every $\alpha$ in the range:
\begin{equation}\label{rangeAl}
0< \alpha <\min\Big\{\frac{\beta}{2},\frac{1}{1+2\gamma}\Big\},
\end{equation}
there exist $\tilde v\in\mathcal{C}^{1,\alpha}(\bar\omega,\R^k)$ and
$\tilde w\in\mathcal{C}^{1,\alpha}(\bar\omega,\R^d)$ with the following properties:
\begin{align*}
& \|\tilde v - v\|_1\leq C\|\mathcal{D}\|_0^{1/2}, \quad \|\tilde w -
w\|_1\leq C\|\mathcal{D}\|_0^{1/2}(1+\|\nabla v\|_0),
\tag*{(\theequation)$_1$}\refstepcounter{equation} \label{Hbound1}\vspace{1mm}\\
& A-\big(\frac{1}{2}(\nabla \tilde v)^T\nabla \tilde v + \sym\nabla
\tilde w\big) =0 \quad\mbox{ in }\; \bar\omega, \tag*{(\theequation)$_2$}\label{Hbound2} 
\end{align*}
%\end{equation}
where the constants $C$ depend only on $d, k$ and $\omega$.
\end{theorem}

\begin{proof}
{\bf 1.} Because of the assumption (\ref{rangeAl}), there exists an exponent:
\begin{equation}\label{exp_del}
\frac{2\gamma\alpha}{(1-\alpha)}< \delta < \min\Big\{1,\frac{2\gamma\beta}{(2-\beta)}\Big\}.
\end{equation} 
We let $\sigma >\sigma_0$ be a sufficiently large constant, in function of
$\delta,\alpha,\gamma$, $\|\nabla v\|_0$ and 
all constants $C$ in the assumed assertions of Theorem
\ref{thm_stage} (these constants depend only in $d,k, \omega$). 

We further set $v_0=v$, $w_0=w$, $\mathcal{D}_0=\mathcal{D}$, and take $M_0\geq \max\{\|v_0\|_2,
\|w_0\|_2,1\}$ that is again sufficiently large, now in function of
$\|A\|_{0,\beta}\|\mathcal{D}\|_0^{\beta/2-1} \sigma^\delta$ and
constants $C$ indicated before.
By successive applications of Theorem \ref{thm_stage} with the chosen $\sigma$ 
and constants $\{M_i\geq 1\}_{i=1}^\infty$ in:
$$M_i = \Big(\tilde C (1+\|\nabla v\|_0) \sigma^\gamma\Big)^iM_0$$
where $\tilde C>1$ is again some large constant (in function of
the aforementioned $C$),
we obtain sequences $\{v_i\in \mathcal{C}^2(\bar\omega,\R^k)\}_{i=1}^\infty$, 
$\{w_i\in\mathcal{C}^2(\bar\omega,\R^d)\}_{i=1}^\infty$ and the related
deficits $\{\mathcal{D}_i \in \mathcal{C}^0(\bar\omega,\R^{d\times d}_\sym)\}_{i=1}^\infty$:
$$\mathcal{D}_i = A-\big(\frac{1}{2}(\nabla v_i)^T\nabla v_i + \sym\nabla w_i\big).$$
We see that, as long as there holds:
\begin{equation}\label{indu_assuD}
0<\|\mathcal{D}_i\|_0\leq 1 \quad \mbox{ and }\quad M_i\geq \max\{\|v_i\|_2,\|w_i\|_2,1\},
\end{equation}
we have, with the constants $C$ depending only on $d, k$ and $\omega$:
\begin{align*}
& \|v_{i+1} - v_i\|_1\leq C\|\mathcal{D}_i\|_0^{1/2}, \quad \|w_{i+1} -
w\|_1\leq C\|\mathcal{D}_i\|_0^{1/2}(1+\|\nabla v_i\|_0),
\tag*{(\theequation)$_1$}\refstepcounter{equation} \label{Bbound1}\vspace{1mm}\\
& \|v_{i+1}\|_2\leq CM_i\sigma^\gamma,\quad \|w_{i+1}\|_2\leq CM_i\sigma^\gamma(1+\|\nabla v_i\|_0),
\tag*{(\theequation)$_2$}\label{Bbound2} \\ 
& \|{\mathcal{D}}_{i+1}\|_0\leq C\Big(\frac{\|A\|_{0,\beta}}{M_i^\beta} \|\mathcal{D}_i\|_0^{\beta/2} +
\frac{\|\mathcal{D}_i\|_0}{\sigma}\Big). \tag*{(\theequation)$_3$} \label{Bbound3}
\end{align*}
%\end{equation}
Below, we inductively validate  (\ref{indu_assuD}) for all $i\geq 0$,
and in fact we show that:
\begin{equation}\label{import_bd}
\|\mathcal{D}_i\|_0\leq \frac{1}{\sigma^{\delta i}}\|\mathcal{D}\|_0\quad
\mbox{ for all } \; i=0\ldots\infty.
\end{equation}

\smallskip

Before doing so, note that (\ref{import_bd}) actually implies both
statements in (\ref{indu_assuD}). Indeed, by \ref{Bbound2}:
\begin{equation*}
\begin{split}
& \|v_{i+1}\|_2\leq CM_i\sigma^\gamma\leq M_{i+1},\vspace{1mm}\\
& \|w_{i+1}\|_2\leq CM_i\sigma^\gamma(1+\|\nabla v_i\|_0) 
\leq CM_i\sigma^\gamma(1+2C +\|\nabla v\|_0) \leq M_{i+1},
\end{split}
\end{equation*}
since by the first bound in \ref{Bbound1} and (\ref{import_bd}) there
follows, provided that $\sigma^{\delta/2}\geq 2$:
\begin{equation}\label{pomoc_v}
\begin{split}
\|\nabla v_i\|_0 & \leq \|\nabla v\|_0 + \sum_{j=0}^{i-1}\|\nabla v_{j+1}-\nabla v_j\|_0
\leq \|\nabla v\|_0 + C\sum_{j=0}^{i-1}\|\mathcal{D}_j\|_0^{1/2} \\ & \leq 
\|\nabla v\|_0 + C\sum_{j=0}^{\infty}\frac{\|\mathcal{D}\|_0^{1/2}}{\sigma^{\delta j/2}} 
= \|\nabla v\|_0 + \frac{C}{1-\sigma^{-\delta/2}} \|\mathcal{D}\|_0^{1/2}
\\ & \leq \|\nabla v\|_0 + {2C} \|\mathcal{D}\|_0^{1/2} \leq 2C +\|\nabla v\|_0.
\end{split}
\end{equation}

\medskip

{\bf 2.} Clearly (\ref{import_bd}) holds at $i=0$. To prove
it at $(i+1)$, use \ref{Bbound3} and the induction assumption:
\begin{equation}\label{dd}
\|\mathcal{D}_{i+1}\|_0\leq C\Big( \frac{\|A\|_{0,\beta}
  \|\mathcal{D}\|_0^{\beta/2}}{M_i^\beta\sigma^{\delta i\beta/2}} +
\frac{\|\mathcal{D}\|_0}{\sigma^{\delta i +1}}\Big) = \frac{\|\mathcal{D}\|_{0}}{\sigma^{\delta(i+1)}}
\Big(\frac{C\|A\|_{0,\beta} \|\mathcal{D}\|_0^{\beta/2-1}}{M_i^\beta\sigma^{\delta i\beta/2 - \delta(i+1)}} +
\frac{C}{\sigma^{1-\delta}}\Big),
\end{equation}
and check that both terms in parentheses in the right hand side above
are not greater than $1/2$. For the second term, this is readily implied by
taking $\sigma$ large enough that $\sigma^{1-\delta}\geq 2C$, in view
of $1-\delta>0$ in (\ref{exp_del}). For the first term, we note that
$$\frac{C\|A\|_{0,\beta} \|\mathcal{D}\|_0^{\beta/2-1}}{M_i^\beta\sigma^{\delta i\beta/2 -  \delta(i+1)}} 
\leq \frac{C\|A\|_{0,\beta} \|\mathcal{D}\|_0^{\beta/2-1}\sigma^\delta}{M_0^\beta}
\cdot \sigma^{\delta i-\gamma\beta i - \delta\beta i/2} \leq \frac{C\|A\|_{0,\beta}
    \|\mathcal{D}\|_0^{\beta/2-1}\sigma^\delta}{M_0^\beta},$$  
since the exponent $\delta i-\gamma\beta i - \delta\beta i/2$ is
non-positive, due to $\delta < \frac{2\gamma\beta}{2-\beta}$ in (\ref{exp_del}):
$$\delta i-\gamma\beta i - \delta\beta i/2 =\frac{i}{2}
\big(\delta(2 - \beta) - 2\gamma\beta \big) \leq 0 \quad\mbox{ for all }\; i\geq 0.$$
In conclusion, the expression in parentheses in (\ref{dd}) is bounded
by $1$, provided $M_0$ has been chosen sufficiently large. This ends the proof of (\ref{import_bd}).

\medskip

{\bf 3.} From \ref{Bbound1}, (\ref{pomoc_v}) and (\ref{import_bd}), it follows that for all $i=0\ldots \infty$:
\begin{equation*}
\begin{split}
& \|v_{i+1} - v_i\|_1\leq \frac{C}{\sigma^{\delta
    i/2}}\|\mathcal{D}\|_0^{1/2}, \qquad \|w_{i+1} - w_i\|_1\leq \frac{C}{\sigma^{\delta
    i/2}}\|\mathcal{D}\|_0^{1/2} \big(1+ \|\nabla v\|_0\big),
\end{split}
\end{equation*}
Hence, both sequences $\{v_i\}_{i=1}^\infty$, 
$\{w_i\}_{i=1}^\infty$ are Cauchy in $\mathcal{C}^1(\bar\omega)$ and as
such they converge to the limit fields, respectively: 
$$\tilde v\in\mathcal{C}^1(\omega,\R^k), \qquad \tilde w\in\mathcal{C}^1(\omega,\R^d)$$ 
that satisfy \ref{Hbound1} and \ref{Hbound2}, in virtue of
$\|\mathcal{D}_i\|_0\to 0$ as $i\to\infty$. 

\smallskip

It remains to show that $\tilde v$ and $\tilde w$ are
$\mathcal{C}^{1,\alpha}$-regular. To this end, we use the estimate:
\begin{equation*}
\begin{split}
\|v_{i+1} - v_i\|_2 + \|w_{i+1} - w_i\|_2 & \leq {C}M_i{\sigma^\gamma}
\big(1+ \|\nabla v\|_0\big) \\ & \leq C \Big(\tilde C (1+\|\nabla
v\|_0)\sigma^\gamma\Big)^{i+1} M_0, 
\end{split}
\end{equation*}
resulting from \ref{Bbound2} and (\ref{pomoc_v}), in the
interpolation inequality $\|\cdot\|_{1,\alpha}\leq C \|\cdot\|_{1}^{1-\alpha}\|\cdot\|_{2}^\alpha$:
\begin{equation*}
\begin{split}
\|v_{i+1} -& v_i\|_{1,\alpha}  + \|w_{i+1} - w_i\|_{1,\alpha} \\ & \leq {C}
\|\mathcal{D}\|_0^{(1-\alpha)/2} \Big(\tilde C (1+\|\nabla
v\|_0)\Big)^{(i+1)\alpha+(1-\alpha)} M_0^\alpha \sigma^{\alpha\gamma(i+1) - \delta i (1-\alpha)/2}
\\ & = {C} \tilde C \cdot M_0^\alpha\|\mathcal{D}\|_0^{(1-\alpha)/2} \big(1+\|\nabla
v\|_0\big) \sigma^{\alpha\gamma} \cdot \Big(\frac{\tilde C^\alpha (1+\|\nabla v\|_0)^{\alpha}}{
\sigma^{\delta (1-\alpha)/2-\alpha\gamma}}\Big)^i.
\end{split}
\end{equation*}
Since the exponent $\delta (1-\alpha)/2-\alpha\gamma$ is positive, in
view of $\delta>\frac{2\gamma\alpha}{1-\alpha}$ in (\ref{exp_del}), we
see that both sequences $\{v_i\}_{i=1}^\infty$, $\{w_i\}_{i=1}^\infty$ are Cauchy in
$\mathcal{C}^{1,\alpha}(\bar\omega)$, provided that $\sigma$ is
sufficiently large to have:
$$\frac{\tilde C^\alpha (1+\|\nabla v\|_0)^{\alpha}}{
\sigma^{\delta (1-\alpha)/2-\alpha\gamma}} <1.$$
In conclusion, $\tilde v\in\mathcal{C}^{1,\alpha}(\omega,\R^k)$ and
$\tilde w\in\mathcal{C}^{1,\alpha}(\omega,\R^d)$ as claimed. The proof
is done.
\end{proof}

\end{document}
